% =============================================================================
\section{Tarefa 4 --- Precisão da cinemática inversa do robô}
\label{sec:ik}
% =============================================================================

Escolheram-se 60 pontos igualmente espaçados no tempo ao longo do swing
(\texttt{src/ik\_compare.py}). Para cada pose pedida $(p_k, R_k)$ obteve-se
$q_k$ com a inversa do controlador\linebreak (\texttt{get\_inverse\_kinematics}),
calculou-se $(\hat p_k, \hat R_k) = \T{0}{6}(q_k)$ com a nossa cinemática
direta e mediram-se o erro de posição $\lVert\hat p_k - p_k\rVert$ e o de
orientação (ângulo de $\hat R_k R_k^{T}$). Para separar o erro da inversa do
erro do nosso modelo, repetiu-se com a cinemática direta do próprio
controlador (\texttt{get\_forward\_kinematics}). O robô não se move neste
ensaio.

\begin{table}[H]
\centering
\caption{Pose recuperada a partir das juntas da inversa do robô, comparada com
a pose pedida (60 pontos).}
\label{tab:ik}
\small
\begin{tabular}{lcccc}
\toprule
\rowcolor{tabcab}
 & \multicolumn{2}{c}{Posição} & \multicolumn{2}{c}{Orientação} \\
\rowcolor{tabcab}
Cinemática direta usada & média & máx. & média & máx. \\
\midrule
Do controlador & \SI{0.1}{\micro\meter} & \SI{1.9}{\micro\meter} & --- & \ang{0.0004} \\
A nossa (Tabela~\ref{tab:dh}) & \SI{1.36}{\milli\meter} & \SI{2.02}{\milli\meter} & \ang{0.014} & \ang{0.030} \\
\bottomrule
\end{tabular}
\end{table}

\begin{figure}[H]
\centering
\fig{Photos/ik_erros.png}{0.66\textwidth}
\caption{Erro da pose recuperada ao longo do swing: posição com as duas
cinemáticas diretas, componentes da diferença com a nossa, e orientação. Os
círculos marcam os pontos em que a inversa devolveu o pulso virado.}
\label{fig:ik}
\end{figure}

\begin{itemize}[itemsep=3pt]
  \item \textbf{A inversa do robô é exata face ao seu modelo.} Resolveu os 60
        pontos, e a direta do controlador devolve a pose pedida com menos de
        \SI{2}{\micro\meter} e \ang{0.0004}, ao nível do arredondamento.
  \item \textbf{A diferença de \SI{1.4}{\milli\meter} está no modelo
        geométrico, não na inversa.} É, ponto a ponto, a diferença entre a
        nossa direta e a do controlador, com um desvio quase constante de
        $(-0{,}2;\,+0{,}6;\,-0{,}9)$~mm. Um erro de \ang{0.03} num braço de
        \SI{300}{\milli\meter} só desloca a ponta \SI{0.16}{\milli\meter}, por
        isso a diferença vem sobretudo de comprimentos. Ajustando pequenas
        correções aos parâmetros da Tabela~\ref{tab:dh}, desce para
        \SI{0.06}{\milli\meter} nestes pontos e de 1,6 para
        \SI{0.9}{\milli\meter} num segundo swing executado no robô, que não
        entrou no ajuste
        (\texttt{src/ajuste\_dh.py}): o controlador usa um modelo que não é
        exatamente o nominal. Com uma só trajetória, porém, vários parâmetros
        compensam-se, e o ajuste não é uma calibração.
        
  \item \textbf{A inversa não garante continuidade das juntas.} Em 19 dos 60
        pontos devolveu a solução de pulso virado
        ($q_4 \pm \ang{180}$, $-q_5$, $q_6 \pm \ang{180}$), a mesma pose com
        outros ângulos; depois de convertidos, todos coincidem com a nossa
        inversa numérica a menos de \ang{0.5}.
\end{itemize}

\begin{destaque}
\textbf{Conclusão.} A inversa incorporada é numericamente exata. A precisão
com que a ferramenta atinge uma pose é limitada pelo modelo geométrico e pela
repetibilidade do robô (1 a \SI{2}{\milli\meter}), e não pela inversa; para
executar trajetórias com as juntas obtidas ponto a ponto é preciso escolher o
ramo da solução.
\end{destaque}
